% !TeX spellcheck = es_ES
\documentclass[12pt]{article}
\usepackage{amsmath}
\usepackage{graphicx}
\usepackage{moodle}
\begin{document}

%Se consideran las siguientes matrices 
%\[
%A=\begin{pmatrix}
%1 & 2 \\ 
%2 & 0
%\end{pmatrix}, 
%\quad 
%B =
%\begin{pmatrix}
%0 & 2 \\ 
%1 & -1
%\end{pmatrix}.
%\]

\begin{quiz}{Matemáticas/Matrices}

 
 	\begin{multi}[shuffle=true,numbering=none]%
		{Matrices 1}
	 	El determinante de $\begin{pmatrix}
	 	1 & 2 \\ 
	 	2 & 0
	 	\end{pmatrix}$ vale
		\item* $-4$ 
		\item $4$
		\item $5$
 	\end{multi}
 
 	\begin{multi}[shuffle=true,numbering=none]%
		{Matrices 2}
	 	El determinante de $\begin{pmatrix}
	 	0 & 2 \\ 
	 	1 & -1
	 	\end{pmatrix}$ vale
		\item $-4$ 
		\item* $-2$
		\item $5$
 	\end{multi}	 
	
 	\begin{multi}[shuffle=true,numbering=none]%
		{Matrices 3}
	 	La inversa de $\begin{pmatrix}
	 	1 & 2 \\ 
	 	2 & 0
	 	\end{pmatrix}$ es
		\item* $\begin{pmatrix} 0 & 1/2 \\ 1/2 & -1/4 \end{pmatrix}$
		\item $\begin{pmatrix} 0 & 1/4 \\ 1/2 & -1/4 \end{pmatrix}$	
		\item $\begin{pmatrix} 0 & 1/2 \\ -1/2 & -1/4 \end{pmatrix}$
		\item $\begin{pmatrix} 1 & 1/2 \\ 1/2 & -1/4 \end{pmatrix}$
 	\end{multi}	

 	\begin{multi}[shuffle=true,numbering=none]%
		{Matrices 4}
	 	El determinante de $\begin{pmatrix}
	 		1 & 2 \\ 
	 		2 & 0
	 	\end{pmatrix}\cdot
	 	\begin{pmatrix}
	 		0 & 2 \\ 
	 		1 & -1
	 	\end{pmatrix}$ vale
		\item* $8$ 
		\item $-4$
		\item $-8$
 	\end{multi}	 

	 			
\end{quiz}
\end{document}
 
